% Practice list A --- topics of Listas 1-2 (Kurlat chs. 1-5), GRADUATE LEVEL (style guide §8).
% \solucoesfalse -> blank list ; \solucoestrue -> answer key + rubric.
% Compile: pdflatex (2x); check: pdffonts (all Type 1). Every number: practice-L1-L2-check.py
\documentclass[11pt,a4paper]{article}
% Shared preamble of the practice lists (practice-L*.tex), adapted from lista-exam.tex.
% Each list sets \solucoesfalse / \solucoestrue and \listatema right after % Shared preamble of the practice lists (practice-L*.tex), adapted from lista-exam.tex.
% Each list sets \solucoesfalse / \solucoestrue and \listatema right after % Shared preamble of the practice lists (practice-L*.tex), adapted from lista-exam.tex.
% Each list sets \solucoesfalse / \solucoestrue and \listatema right after \input{practice-preamble}.
\usepackage[utf8]{inputenc}
\usepackage[T1]{fontenc}
\usepackage{lmodern}   % Type 1 vector fonts (replaces PK bitmaps)
\usepackage{cmap}      % ToUnicode CMap -> copyable text and accents
\usepackage[english]{babel}
\usepackage{amsmath,amssymb,amsfonts}
\usepackage[a4paper,margin=2.2cm]{geometry}
\usepackage{enumitem}
\usepackage{fancyhdr}
\usepackage{comment}
\usepackage{xcolor}
\usepackage{microtype}
\usepackage{booktabs}
\usepackage{tikz}

\newif\ifsolucoes
\newcommand{\disciplina}{Macroeconomics I --- MPE Insper}
\newcommand{\espacoresol}{4.5cm}   % height of the answer space in the blank list
% The sol environment must be decided after the list sets its toggle.
\AtBeginDocument{\ifsolucoes\includecomment{sol}\else\excludecomment{sol}\fi}

\pagestyle{fancy}\fancyhf{}
\lhead{\footnotesize\disciplina}
\rhead{\footnotesize Practice list\ifsolucoes{} --- Solutions and rubric\fi}
\cfoot{\footnotesize\thepage}
\renewcommand{\headrulewidth}{0.3pt}
\setlength{\headheight}{14pt}
\setlist{topsep=2pt,itemsep=2pt}
\newcounter{questao}

\newcommand{\bloco}[3]{\bigskip\noindent
  {\large\bfseries\color{blue!55!black}Block #1 --- #2}\par\noindent\hfill\textit{(#3)}\par
  \noindent\rule{\linewidth}{0.3pt}\smallskip}
\newcommand{\questao}[1]{\refstepcounter{questao}\medskip\noindent
  \textbf{Question \thequestao.}~#1\par}
\newenvironment{alternativas}{\begin{enumerate}[label=\Alph*.,leftmargin=2.2em,topsep=1pt]}{\end{enumerate}}
\newcommand{\gabarito}[1]{\ifsolucoes\par\noindent\textbf{Answer:}~#1\fi}
\newcommand{\resolespaco}{\ifsolucoes\else\par\textit{Solution:}\vspace{\espacoresol}\par\fi}
\newcommand{\passo}[2]{\par\smallskip\noindent\textbf{Step --- #1.}~#2}
\newcommand{\intuicao}[1]{\par\smallskip\noindent\textit{Economic intuition.}~#1}
\newcommand{\refbib}[1]{\par\smallskip\noindent{\footnotesize\color{black!60}\textbf{Ref:} #1}}
\newcommand{\rubrica}[1]{\par\smallskip\noindent\textbf{Grading rubric.}\begin{itemize}[leftmargin=1.4em,topsep=1pt]#1\end{itemize}}

\newcommand{\cabecalho}{%
\begin{center}
  {\Large\bfseries \disciplina}\\[2pt]
  {\large Practice list: \listatema}\ifsolucoes\\[2pt]\textit{Solutions and grading rubric}\fi
\end{center}
\vspace{2pt}
{\footnotesize\noindent\textbf{Instructions.} Multiple choice: mark the single correct option (no
penalty for wrong answers). \emph{True/False}: every 2 wrong items cancel 1 right one
(score $=\max\{0,\text{right}-\text{wrong}/2\}\times$ item value); blank items do not count.
Discursive questions: show every step. Numbers are illustrative and chosen to give clean answers.\par}}
.
\usepackage[utf8]{inputenc}
\usepackage[T1]{fontenc}
\usepackage{lmodern}   % Type 1 vector fonts (replaces PK bitmaps)
\usepackage{cmap}      % ToUnicode CMap -> copyable text and accents
\usepackage[english]{babel}
\usepackage{amsmath,amssymb,amsfonts}
\usepackage[a4paper,margin=2.2cm]{geometry}
\usepackage{enumitem}
\usepackage{fancyhdr}
\usepackage{comment}
\usepackage{xcolor}
\usepackage{microtype}
\usepackage{booktabs}
\usepackage{tikz}

\newif\ifsolucoes
\newcommand{\disciplina}{Macroeconomics I --- MPE Insper}
\newcommand{\espacoresol}{4.5cm}   % height of the answer space in the blank list
% The sol environment must be decided after the list sets its toggle.
\AtBeginDocument{\ifsolucoes\includecomment{sol}\else\excludecomment{sol}\fi}

\pagestyle{fancy}\fancyhf{}
\lhead{\footnotesize\disciplina}
\rhead{\footnotesize Practice list\ifsolucoes{} --- Solutions and rubric\fi}
\cfoot{\footnotesize\thepage}
\renewcommand{\headrulewidth}{0.3pt}
\setlength{\headheight}{14pt}
\setlist{topsep=2pt,itemsep=2pt}
\newcounter{questao}

\newcommand{\bloco}[3]{\bigskip\noindent
  {\large\bfseries\color{blue!55!black}Block #1 --- #2}\par\noindent\hfill\textit{(#3)}\par
  \noindent\rule{\linewidth}{0.3pt}\smallskip}
\newcommand{\questao}[1]{\refstepcounter{questao}\medskip\noindent
  \textbf{Question \thequestao.}~#1\par}
\newenvironment{alternativas}{\begin{enumerate}[label=\Alph*.,leftmargin=2.2em,topsep=1pt]}{\end{enumerate}}
\newcommand{\gabarito}[1]{\ifsolucoes\par\noindent\textbf{Answer:}~#1\fi}
\newcommand{\resolespaco}{\ifsolucoes\else\par\textit{Solution:}\vspace{\espacoresol}\par\fi}
\newcommand{\passo}[2]{\par\smallskip\noindent\textbf{Step --- #1.}~#2}
\newcommand{\intuicao}[1]{\par\smallskip\noindent\textit{Economic intuition.}~#1}
\newcommand{\refbib}[1]{\par\smallskip\noindent{\footnotesize\color{black!60}\textbf{Ref:} #1}}
\newcommand{\rubrica}[1]{\par\smallskip\noindent\textbf{Grading rubric.}\begin{itemize}[leftmargin=1.4em,topsep=1pt]#1\end{itemize}}

\newcommand{\cabecalho}{%
\begin{center}
  {\Large\bfseries \disciplina}\\[2pt]
  {\large Practice list: \listatema}\ifsolucoes\\[2pt]\textit{Solutions and grading rubric}\fi
\end{center}
\vspace{2pt}
{\footnotesize\noindent\textbf{Instructions.} Multiple choice: mark the single correct option (no
penalty for wrong answers). \emph{True/False}: every 2 wrong items cancel 1 right one
(score $=\max\{0,\text{right}-\text{wrong}/2\}\times$ item value); blank items do not count.
Discursive questions: show every step. Numbers are illustrative and chosen to give clean answers.\par}}
.
\usepackage[utf8]{inputenc}
\usepackage[T1]{fontenc}
\usepackage{lmodern}   % Type 1 vector fonts (replaces PK bitmaps)
\usepackage{cmap}      % ToUnicode CMap -> copyable text and accents
\usepackage[english]{babel}
\usepackage{amsmath,amssymb,amsfonts}
\usepackage[a4paper,margin=2.2cm]{geometry}
\usepackage{enumitem}
\usepackage{fancyhdr}
\usepackage{comment}
\usepackage{xcolor}
\usepackage{microtype}
\usepackage{booktabs}
\usepackage{tikz}

\newif\ifsolucoes
\newcommand{\disciplina}{Macroeconomics I --- MPE Insper}
\newcommand{\espacoresol}{4.5cm}   % height of the answer space in the blank list
% The sol environment must be decided after the list sets its toggle.
\AtBeginDocument{\ifsolucoes\includecomment{sol}\else\excludecomment{sol}\fi}

\pagestyle{fancy}\fancyhf{}
\lhead{\footnotesize\disciplina}
\rhead{\footnotesize Practice list\ifsolucoes{} --- Solutions and rubric\fi}
\cfoot{\footnotesize\thepage}
\renewcommand{\headrulewidth}{0.3pt}
\setlength{\headheight}{14pt}
\setlist{topsep=2pt,itemsep=2pt}
\newcounter{questao}

\newcommand{\bloco}[3]{\bigskip\noindent
  {\large\bfseries\color{blue!55!black}Block #1 --- #2}\par\noindent\hfill\textit{(#3)}\par
  \noindent\rule{\linewidth}{0.3pt}\smallskip}
\newcommand{\questao}[1]{\refstepcounter{questao}\medskip\noindent
  \textbf{Question \thequestao.}~#1\par}
\newenvironment{alternativas}{\begin{enumerate}[label=\Alph*.,leftmargin=2.2em,topsep=1pt]}{\end{enumerate}}
\newcommand{\gabarito}[1]{\ifsolucoes\par\noindent\textbf{Answer:}~#1\fi}
\newcommand{\resolespaco}{\ifsolucoes\else\par\textit{Solution:}\vspace{\espacoresol}\par\fi}
\newcommand{\passo}[2]{\par\smallskip\noindent\textbf{Step --- #1.}~#2}
\newcommand{\intuicao}[1]{\par\smallskip\noindent\textit{Economic intuition.}~#1}
\newcommand{\refbib}[1]{\par\smallskip\noindent{\footnotesize\color{black!60}\textbf{Ref:} #1}}
\newcommand{\rubrica}[1]{\par\smallskip\noindent\textbf{Grading rubric.}\begin{itemize}[leftmargin=1.4em,topsep=1pt]#1\end{itemize}}

\newcommand{\cabecalho}{%
\begin{center}
  {\Large\bfseries \disciplina}\\[2pt]
  {\large Practice list: \listatema}\ifsolucoes\\[2pt]\textit{Solutions and grading rubric}\fi
\end{center}
\vspace{2pt}
{\footnotesize\noindent\textbf{Instructions.} Multiple choice: mark the single correct option (no
penalty for wrong answers). \emph{True/False}: every 2 wrong items cancel 1 right one
(score $=\max\{0,\text{right}-\text{wrong}/2\}\times$ item value); blank items do not count.
Discursive questions: show every step. Numbers are illustrative and chosen to give clean answers.\par}}

\solucoesfalse                 % <<<< \solucoestrue for the ANSWER KEY
\newcommand{\listatema}{Measurement, Welfare and Growth --- graduate level (Listas 1--2)}
\renewcommand{\espacoresol}{6cm}

\begin{document}
\cabecalho

% =============================================================================
\bloco{1}{Price indices, PPP and welfare across countries}{Kurlat chs.~1--2 \,\textbullet\, 25 points}

\questao{(10 points) Objective items. Each item needs a short derivation; write it on scratch paper.}

\textbf{a)} (3 points) A household has preferences $u(x_1,x_2)=x_1^{1/3}x_2^{2/3}$ over electricity
($x_1$) and everything else ($x_2$). Between two months the price of electricity rises 44\% and the
other price does not change. The true cost-of-living (Konüs) index, defined as the ratio of minimum
expenditures needed to reach the \emph{same} utility at the new and old prices, records inflation of:
\begin{alternativas}
\item $14.67\%$
\item $12.92\%$
\item $12.99\%$
\item $11.34\%$
\end{alternativas}
\gabarito{B}
\begin{sol}
\passo{expenditure function}{$\min p_1x_1+p_2x_2$ s.t.\ $x_1^{1/3}x_2^{2/3}\ge\bar u$. The FOCs give
$x_1=\tfrac13 e/p_1$, $x_2=\tfrac23 e/p_2$ (constant expenditure shares $s_1=\tfrac13$, $s_2=\tfrac23$).
Substituting into the constraint: $e(p,\bar u)=\bar u\,(3p_1)^{1/3}(\tfrac32p_2)^{2/3}$.}
\passo{Konüs index}{$P^*=\dfrac{e(p^1,\bar u)}{e(p^0,\bar u)}=\Big(\dfrac{p_1^1}{p_1^0}\Big)^{1/3}\Big(\dfrac{p_2^1}{p_2^0}\Big)^{2/3}=1.44^{1/3}=1.1292$:
the geometric mean of the price relatives, weighted by the shares, independent of $\bar u$ (homotheticity).}
\passo{distractors}{With constant shares, Laspeyres $=\sum_is_iR_i=\tfrac13 1.44+\tfrac23=1.1467$ (A);
Paasche $=(\sum_is_i/R_i)^{-1}=1.1134$ (D); Fisher $=\sqrt{1.1467\cdot1.1134}=1.1299$ (C).}
\refbib{Kurlat (2020), \S1.2, pp.~22--25; Jones (2020), \S2.3, pp.~32--33.}
\end{sol}

\medskip
\textbf{b)} (3 points) Brazil and the US produce a tradable good with $Y_T=A_TL_T$ and a non-tradable
service with $Y_N=A_NL_N$. Labor is mobile across sectors within each country, markets are
competitive, the tradable obeys the law of one price, and consumers' price index is
$P=p_T^{1/2}p_N^{1/2}$. Brazil's productivity relative to the US is $0.25$ in tradables and $0.80$ in
non-tradables. Brazil's price level relative to the US, both measured in dollars at the market
exchange rate, is:
\begin{alternativas}
\item $0.250$
\item $0.800$
\item $0.559$
\item $0.313$
\end{alternativas}
\gabarito{C}
\begin{sol}
\passo{zero profits}{With a common wage $W$: $p_T=W/A_T$, $p_N=W/A_N$, so $p_N/p_T=A_T/A_N$.}
\passo{price level}{$P=p_T(p_N/p_T)^{1/2}=p_T(A_T/A_N)^{1/2}$. The law of one price sets
$p_T^{BR}=E\,p_T^{US}$, hence
$\dfrac{P^{BR}}{E\,P^{US}}=\Big[\dfrac{(A_T/A_N)^{BR}}{(A_T/A_N)^{US}}\Big]^{1/2}=\Big(\dfrac{0.25}{0.80}\Big)^{1/2}=0.559$.
Option D forgets the share $1/2$.}
\intuicao{Low tradable productivity means low wages; low wages make services cheap. Market-rate
GDP therefore understates Brazil's real output by a factor $1/0.559\approx1.79$ (for a GDP basket
like the consumption basket).}
\refbib{Kurlat (2020), \S1.2, pp.~24--27; Romer (2012), \S4.2, eq.\ (4.16), pp.~160--161.}
\end{sol}

\medskip
\textbf{c)} (2 points, T or F) With Cobb--Douglas preferences the Fisher index equals the true
cost-of-living index, since it is the geometric mean of an upper bound (Laspeyres) and a lower
bound (Paasche).
\gabarito{F}
\begin{sol}
\passo{three means}{With constant shares, Laspeyres is the arithmetic mean (AM), Paasche the harmonic
mean (HM) and the true index the geometric mean (GM) of the price relatives. HM $\le$ GM $\le$ AM is
true, but $\sqrt{\text{AM}\cdot\text{HM}}=\text{GM}$ only in special cases (two goods with equal
shares, where AM$\cdot$HM$=R_1R_2$). In (a): Fisher $12.99\%\ne12.92\%$.}
\refbib{Kurlat (2020), \S1.2.}
\end{sol}

\medskip
\textbf{d)} (2 points, T or F) If log consumption is normal with variance $v$ and utility is CRRA with
curvature $\sigma$, the consumption-equivalent cost of inequality, $1-e^{-\sigma v/2}$, depends
neither on mean consumption nor on the value of life $\bar u$, while the welfare cost of shorter
life expectancy depends on both.
\gabarito{T}
\begin{sol}
\passo{why}{Question 2(a) proves $c_{CE}=\bar c\,e^{-\sigma v/2}$: a proportional factor, free of $\bar c$
and $\bar u$. Life expectancy multiplies the whole flow utility $\bar u+u(c_{CE})-\theta(1-l)^2$,
which depends on both (Question 2(b)).}
\refbib{Kurlat (2020), \S2.2, eqs.\ (2.2.1)--(2.2.5).}
\end{sol}

% -----------------------------------------------------------------------------
\questao{(15 points) Discursive. The Jones--Klenow measure (Kurlat \S2.2). ``Rawls'' is placed at
random in country $j\in\{BR,US\}$. He is alive with probability $a_j$ (life expectancy over 100) and,
if alive, gets $\bar u+u(c)-\theta(1-l_j)^2$, where $1-l_j$ is market work (the same for everyone) and
$u(c)=\frac{c^{1-\sigma}-1}{1-\sigma}$ ($\ln c$ if $\sigma=1$). Consumption is lognormal:
$\ln c\sim N(m_j,v_j)$, with mean $\bar c_j=e^{m_j+v_j/2}$. The welfare of Brazil relative to the US
is the $\lambda$ that makes Rawls indifferent between Brazil and a US in which \emph{every}
consumption level is multiplied by $\lambda$.
Data: $\bar c_{US}=1$, $\bar c_{BR}=0.30$, $v_{US}=0.4$, $v_{BR}=1.0$, $a_{US}=0.79$, $a_{BR}=0.76$,
$1-l_{US}=0.20$, $1-l_{BR}=0.22$, $\theta=2$, $\bar u=5$.}

\textbf{a)} (3 points) Prove that $\mathbb E[u(c)]=u(c_{CE})$ with $c_{CE}=\bar c\,e^{-\sigma v/2}$ for
every $\sigma>0$, including $\sigma=1$. Interpret the factor $e^{-\sigma v/2}$.
\resolespaco
\begin{sol}
\passo{normal moment generating function}{If $X\sim N(m,v)$ then $\mathbb E[e^{tX}]=e^{tm+t^2v/2}$ for every real $t$.}
\passo{case $\sigma\ne1$}{With $X=\ln c$ and $t=1-\sigma$:
\[\mathbb E[c^{1-\sigma}]=e^{(1-\sigma)m+(1-\sigma)^2v/2}.\]
Substitute $m=\ln\bar c-v/2$:
\[\mathbb E[c^{1-\sigma}]=\bar c^{\,1-\sigma}e^{-(1-\sigma)v/2+(1-\sigma)^2v/2}
=\bar c^{\,1-\sigma}e^{-\sigma(1-\sigma)v/2}=\big(\bar c\,e^{-\sigma v/2}\big)^{1-\sigma}.\]
Since $u$ is affine in $c^{1-\sigma}$, $\mathbb E[u(c)]=u(\bar c\,e^{-\sigma v/2})$.}
\passo{case $\sigma=1$}{$\mathbb E[\ln c]=m=\ln\bar c-v/2=\ln(\bar c\,e^{-v/2})$: the same formula.}
\passo{interpretation}{Lognormal inequality acts as a proportional tax $e^{-\sigma v/2}$ on mean
consumption. It grows with dispersion $v$ and with curvature $\sigma$: behind the veil of ignorance,
inequality is risk, and $\sigma$ is risk aversion. It does not depend on $\bar c$.}
\refbib{Kurlat (2020), \S2.2, eqs.\ (2.2.4)--(2.2.6), pp.~35--37.}
\rubrica{
\item Uses the moment generating function correctly: 1 pt.
\item Algebra to $(\bar c e^{-\sigma v/2})^{1-\sigma}$ and the affine step: 1 pt.
\item Log case and interpretation: 1 pt. Stating the result without proof: max.\ 1 pt in the item.
}
\end{sol}

\textbf{b)} (4 points) Let $\sigma=1$. Derive $\lambda$ and an \emph{exact} additive decomposition of
$\ln\lambda$ into consumption, inequality, leisure and life-expectancy terms. Compute every term.
\resolespaco
\begin{sol}
\passo{expected utility}{Scaling every consumption level by $\lambda$ adds $\ln\lambda$ to $\ln c$.
Using (a), with $h_j\equiv1-l_j$:
\[U_j(\lambda)=a_j\big[\bar u+\ln\lambda+\ln\bar c_j-\tfrac{v_j}{2}-\theta h_j^2\big].\]}
\passo{indifference}{$U_{US}(\lambda)=U_{BR}(1)$. Define $X_{BR}\equiv\bar u+\ln\bar c_{BR}-v_{BR}/2-\theta h_{BR}^2$
(flow utility of an average-risk Brazilian). Then
\[\ln\lambda=\frac{a_{BR}}{a_{US}}X_{BR}-\bar u-\ln\bar c_{US}+\frac{v_{US}}{2}+\theta h_{US}^2.\]}
\passo{exact decomposition}{Add and subtract $X_{BR}$:
\[\ln\lambda=\underbrace{\ln\frac{\bar c_{BR}}{\bar c_{US}}}_{\text{consumption}}
\underbrace{-\frac{v_{BR}-v_{US}}{2}}_{\text{inequality}}
\underbrace{-\theta\big(h_{BR}^2-h_{US}^2\big)}_{\text{leisure}}
+\underbrace{\Big(\frac{a_{BR}}{a_{US}}-1\Big)X_{BR}}_{\text{life expectancy}}.\]
No approximation is involved.}
\passo{numbers}{Consumption $\ln0.30=-1.20397$; inequality $-0.30000$; leisure
$-2(0.0484-0.0400)=-0.01680$; $X_{BR}=5-1.20397-0.5-0.0968=3.19923$; life expectancy
$(0.76/0.79-1)\cdot3.19923=-0.12149$. Total $\ln\lambda=-1.64226$, so $\lambda=0.1936$.}
\intuicao{Brazil's welfare is about 19\% of the US level, against a consumption ratio of 30\%:
inequality removes 0.30 log points, shorter lives 0.12, longer hours 0.02.}
\refbib{Kurlat (2020), \S2.2, eqs.\ (2.2.1)--(2.2.2), pp.~33--40.}
\rubrica{
\item Expected utility with $\mathbb E\ln c=\ln\bar c-v/2$ and $a_j$ multiplying the whole bracket: 1 pt.
\item Solves for $\ln\lambda$: 1 pt.
\item Exact decomposition: 1 pt. A decomposition that leaves an unexplained residual: max.\ 0.5 pt on this line.
\item Numbers: 1 pt ($-1$ pt for an arithmetic slip with the right method).
}
\end{sol}

\textbf{c)} (4 points) Now let $\sigma$ be general. Derive $\lambda$ in closed form. Prove that it tends
to the formula of (b) as $\sigma\to1$. Prove that the pure inequality factor (all else equal across
countries) falls with $\sigma$ when $v_{BR}>v_{US}$, and give its limits as $\sigma\to0$ and
$\sigma\to\infty$. Explain why Kurlat's unnormalized $c^{1-\sigma}/(1-\sigma)$ with a fixed $\bar u$ has
no limit at $\sigma=1$.
\resolespaco
\begin{sol}
\passo{indifference}{By (a), scaling every consumption by $\lambda$ scales $c_{CE}$ by $\lambda$:
$U_{US}(\lambda)=a_{US}[\bar u+u(\lambda c^{US}_{CE})-\theta h_{US}^2]$. Indifference with Brazil:
\[u(\lambda c^{US}_{CE})=Z\equiv\frac{a_{BR}}{a_{US}}\big[\bar u+u(c^{BR}_{CE})-\theta h_{BR}^2\big]-\bar u+\theta h_{US}^2.\]}
\passo{inversion}{$u^{-1}(Z)=[1+(1-\sigma)Z]^{1/(1-\sigma)}$, so
\[\lambda=\frac{[1+(1-\sigma)Z]^{1/(1-\sigma)}}{c^{US}_{CE}},\qquad c^{j}_{CE}=\bar c_je^{-\sigma v_j/2},\]
defined when $1+(1-\sigma)Z>0$.}
\passo{limit $\sigma\to1$}{By L'Hôpital in $\sigma$, $\frac{c^{1-\sigma}-1}{1-\sigma}\to\ln c$, so $Z\to Z_{\log}$; and
$[1+(1-\sigma)Z]^{1/(1-\sigma)}\to e^{Z}$ (the definition of $e$). Hence $\lambda\to e^{Z_{\log}}/c^{US}_{CE}$,
which is exactly the formula of (b).}
\passo{inequality factor}{With equal $a$, $h$ and $\bar c$, $\lambda=c^{BR}_{CE}/c^{US}_{CE}=e^{-\sigma(v_{BR}-v_{US})/2}$, and
$\partial\lambda/\partial\sigma=-\tfrac{v_{BR}-v_{US}}{2}\lambda<0$. As $\sigma\to0$ the factor $\to1$ (risk
neutrality: inequality is irrelevant); as $\sigma\to\infty$ it $\to0$ (a maximin, Rawlsian judgment;
lognormal consumption has no positive lower bound).}
\passo{normalization}{$\frac{c^{1-\sigma}}{1-\sigma}=\frac{c^{1-\sigma}-1}{1-\sigma}+\frac{1}{1-\sigma}$. The
constant $1/(1-\sigma)$ diverges at $\sigma=1$; with $\bar u$ held fixed it changes the value of being
alive by an unbounded amount, so $\lambda$ explodes. The value of life must be recalibrated
whenever $\sigma$ changes; normalizing $u$ is one consistent way.}
\refbib{Kurlat (2020), \S2.2, footnote 6 and Exercise 2.10 (Rawlsian limit).}
\rubrica{
\item Indifference condition and inversion: 1.5 pts.
\item Limit at $\sigma=1$ with an argument (L'Hôpital or series): 1 pt. Asserting it: 0 for this line.
\item Sign of $\partial\lambda/\partial\sigma$ and both limits: 1 pt.
\item Normalization argument: 0.5 pt.
}
\end{sol}

\textbf{d)} (4 points) Compute the inequality factor at $\sigma=1$ and $\sigma=2$. Then show that with
$\sigma=2$ and $\bar u=5$ the flow utility of the Brazilian certainty-equivalent resident is
negative, and explain why that makes $\lambda$ meaningless. What must the analyst do?
\resolespaco
\begin{sol}
\passo{factors}{$e^{-1\cdot0.6/2}=0.7408$ and $e^{-2\cdot0.6/2}=0.5488$: the cost of Brazil's
extra inequality rises from 26\% to 45\% of consumption.}
\passo{negative flow utility}{$c^{BR}_{CE}=0.30\,e^{-1.0}=0.1104$; $u=1-1/0.1104=-8.061$; flow
utility $5-8.061-2(0.22)^2=-3.158<0$.}
\passo{why it breaks}{With negative flow utility, being alive is worse than being dead (who gets
zero). A higher survival probability $a_{BR}$ would \emph{lower} welfare, and $\lambda$ would reward
shorter lives.}
\passo{remedy}{Recalibrate $\bar u$ jointly with $\sigma$ (Jones--Klenow match the value of a
statistical life), so that flow utility is positive over the relevant consumption range.}
\intuicao{$\sigma$ does double duty: it is intertemporal curvature, risk aversion and aversion to
inequality at once. Welfare rankings that move with $\sigma$ should be reported with it.}
\refbib{Kurlat (2020), \S2.2, ``The value of life expectancy'', pp.~38--39; Exercise 2.6.}
\rubrica{
\item Both factors: 1 pt.
\item Negative flow utility computed: 1 pt.
\item The perverse effect of life expectancy explained: 1 pt.
\item Recalibration of $\bar u$: 1 pt.
}
\end{sol}

% =============================================================================
\bloco{2}{Solow: CES, exact dynamics and human capital}{Kurlat chs.~4--5; Romer ch.~1 \,\textbullet\, 25 points}

\questao{(10 points) Objective items. Work \emph{per effective worker}: $\tilde k=K/(AL)$,
$\dot{\tilde k}=s f(\tilde k)-(n+g+\delta)\tilde k$ in continuous time.}

\textbf{a)} (3 points) Let $f(\tilde k)=[\alpha\tilde k^{\rho}+1-\alpha]^{1/\rho}$ with $\alpha=1/3$,
$\rho=1/2$ (elasticity of substitution 2), and $n+g+\delta=0.08$. The saving rate above which
$\tilde k$ grows forever, with no steady state, is:
\begin{alternativas}
\item $24\%$
\item $33.3\%$
\item $72\%$
\item none: the Solow model always has a steady state
\end{alternativas}
\gabarito{C}
\begin{sol}
\passo{growth rate}{$\dot{\tilde k}/\tilde k=s\,f(\tilde k)/\tilde k-(n+g+\delta)$, and
$f(k)/k=[\alpha+(1-\alpha)k^{-\rho}]^{1/\rho}$ falls from $\infty$ ($k\to0$) to $\alpha^{1/\rho}=1/9$ ($k\to\infty$).}
\passo{threshold}{A steady state exists iff $s\alpha^{1/\rho}<n+g+\delta$, i.e.\ $s<0.08\cdot9=72\%$.
Above it, $\tilde k$ grows asymptotically at $s/9-0.08>0$: an AK limit. Option A uses $\alpha$
instead of $\alpha^{1/\rho}$; D assumes the Inada conditions, which fail when $\rho>0$.}
\refbib{Romer (2012), \S1.8, pp.~42--43; Kurlat (2020), Exercise 4.4.}
\end{sol}

\medskip
\textbf{b)} (3 points) Let $f(\tilde k)=\tilde k^{1/3}$ and $n+g+\delta=0.08$. The economy starts at
$\tilde k(0)=\tilde k^*/8$. The time until $\tilde k$ reaches $\tilde k^*/2$ is:
\begin{alternativas}
\item $13.0$ years
\item $20.6$ years
\item $13.2$ years
\item $6.6$ years
\end{alternativas}
\gabarito{C}
\begin{sol}
\passo{exact solution}{Let $z\equiv\tilde k^{1-\alpha}$ ($=K/Y$). Then
$\dot z=(1-\alpha)\tilde k^{-\alpha}\dot{\tilde k}=(1-\alpha)[s-(n+g+\delta)z]$, a linear ODE:
$z(t)=z^*+(z(0)-z^*)e^{-\lambda t}$, $\lambda=(1-\alpha)(n+g+\delta)=0.0533$.}
\passo{timing}{$z/z^*=(\tilde k/\tilde k^*)^{2/3}$: $z(0)/z^*=1/4$. Solve $1-0.75e^{-\lambda t}=0.5^{2/3}=0.6300$:
$e^{-\lambda t}=0.4934$, $t=0.7065/0.0533=13.2$ years.}
\passo{distractors}{B is the log-linear approximation ($\ln(\tilde k/\tilde k^*)$ shrinking at rate $\lambda$:
$e^{-\lambda t}=\ln2/\ln8$, $t=20.6$), poor this far from the steady state; A is the half-life $\ln2/\lambda$.}
\refbib{Romer (2012), \S1.5, eqs.\ (1.27)--(1.31), pp.~24--26.}
\end{sol}

\medskip
\textbf{c)} (2 points, T or F) With CES production and an elasticity of substitution below one
($\rho<0$), the capital share $\alpha_K(\tilde k)=f'(\tilde k)\tilde k/f(\tilde k)$ falls as $\tilde k$ rises.
\gabarito{T}
\begin{sol}
\passo{share}{$\alpha_K=\dfrac{\alpha k^{\rho}}{\alpha k^{\rho}+1-\alpha}=\Big[1+\tfrac{1-\alpha}{\alpha}k^{-\rho}\Big]^{-1}$;
with $\rho<0$, $k^{-\rho}$ rises with $k$, so $\alpha_K$ falls: with complements, the abundant
factor's share shrinks. The convergence speed, $(1-\alpha_K(\tilde k^*))(n+g+\delta)$, rises with $\tilde k^*$.}
\refbib{Romer (2012), \S1.5 and \S1.8.}
\end{sol}

\medskip
\textbf{d)} (2 points, T or F) On a balanced growth path with $Y=K^{\alpha}(AL)^{1-\alpha}$ and $A$ growing at
$g$, growth accounting attributes all of the growth of output per worker to the Solow residual.
\gabarito{F}
\begin{sol}
\passo{accounting}{On the path $g_K=g_Y=n+g$, $g_L=n$: residual $=g_Y-\alpha g_K-(1-\alpha)g_L=(1-\alpha)g$;
capital deepening $\alpha(g_K-g_L)=\alpha g$. Accounting assigns a share $\alpha$ to capital, although
technology causes all of it: capital deepening is induced by $A$.}
\refbib{Kurlat (2020), \S5.4, eq.\ (5.4.2), p.~87; Romer (2012), problem 1.13.}
\end{sol}

% -----------------------------------------------------------------------------
\questao{(15 points) Discursive. Mankiw--Romer--Weil with Brazilian data. Output is
\[Y=K^{\alpha}H^{\beta}(AL)^{1-\alpha-\beta},\qquad\alpha+\beta<1,\]
with $\dot K=s_kY-\delta K$ and $\dot H=s_hY-\delta H$; $L$ and $A$ grow at $n$ and $g$. Let
$k=K/(AL)$, $h=H/(AL)$, $y=k^{\alpha}h^{\beta}$ and $N\equiv n+g+\delta$.}

\textbf{a)} (3 points) Derive the laws of motion of $k$ and $h$, the steady state $(k^*,h^*)$ and
$\ln y^*$. Give the elasticity of $y^*$ with respect to $s_k$.
\resolespaco
\begin{sol}
\passo{laws of motion}{$\dot k/k=\dot K/K-n-g=s_ky/k-N$, so $\dot k=s_kk^{\alpha}h^{\beta}-Nk$; likewise
$\dot h=s_hk^{\alpha}h^{\beta}-Nh$.}
\passo{steady state}{$s_ky=Nk$ and $s_hy=Nh$ give $h/k=s_h/s_k$. Substituting into the first:
$s_kk^{\alpha}(ks_h/s_k)^{\beta}=Nk\Rightarrow k^{1-\alpha-\beta}=s_k^{1-\beta}s_h^{\beta}/N$:
\[k^*=\Big(\frac{s_k^{1-\beta}s_h^{\beta}}{N}\Big)^{\frac1{1-\alpha-\beta}},\qquad
h^*=\Big(\frac{s_k^{\alpha}s_h^{1-\alpha}}{N}\Big)^{\frac1{1-\alpha-\beta}}.\]}
\passo{output}{$\ln y^*=\alpha\ln k^*+\beta\ln h^*=\dfrac{\alpha\ln s_k+\beta\ln s_h-(\alpha+\beta)\ln N}{1-\alpha-\beta}$.
Elasticity to $s_k$: $\alpha/(1-\alpha-\beta)$, which is $1$ with $\alpha=\beta=1/3$ (against $0.5$ in the
textbook Solow model).}
\refbib{Romer (2012), problems 4.5--4.9, pp.~185--186; Kurlat (2020), \S5.2.}
\rubrica{
\item Laws of motion with $N=n+g+\delta$: 1 pt.
\item Steady state: 1 pt.
\item $\ln y^*$ and elasticity: 1 pt.
}
\end{sol}

\textbf{b)} (4 points) Prove local stability. Linearize the system in $x=\ln k$, $z=\ln h$, compute the
Jacobian at the steady state, its eigenvalues, and the speed at which $\ln y$ converges. Compare with
the textbook Solow model ($\beta=0$) for $\alpha=\beta=1/3$, $N=0.06$.
\resolespaco
\begin{sol}
\passo{system in logs}{$\dot x=s_ke^{(\alpha-1)x+\beta z}-N$, $\dot z=s_he^{\alpha x+(\beta-1)z}-N$. At the steady
state each exponential term equals $N$.}
\passo{Jacobian}{Differentiating, $J=N\begin{pmatrix}\alpha-1&\beta\\ \alpha&\beta-1\end{pmatrix}$, with
$\operatorname{tr}J=N(\alpha+\beta-2)<0$ and $\det J=N^2[(\alpha-1)(\beta-1)-\alpha\beta]=N^2(1-\alpha-\beta)>0$.}
\passo{eigenvalues}{Discriminant $(\operatorname{tr}J)^2-4\det J=N^2(\alpha+\beta)^2\ge0$, so the roots are real:
\[\mu_{1,2}=\tfrac N2\big[(\alpha+\beta-2)\pm(\alpha+\beta)\big]\in\{-N(1-\alpha-\beta),\,-N\}.\]
Both are negative: the steady state is a locally stable node.}
\passo{speed}{$\ln y=\alpha x+\beta z$ converges at the slow root, $\lambda=(1-\alpha-\beta)N=2\%$ a year:
half-life $\ln2/0.02=34.7$ years. Textbook Solow: $(1-\alpha)N=4\%$, half-life $17.3$ years.}
\intuicao{A broad capital share $\alpha+\beta=2/3$ weakens diminishing returns, so convergence is slow,
in line with the 2\% estimates of cross-country regressions.}
\refbib{Romer (2012), \S1.5, pp.~24--26, and problem 4.8.}
\rubrica{
\item Log system and Jacobian: 1.5 pts.
\item Trace/determinant argument for stability: 1 pt. Stability asserted without it: 0 for this line.
\item Eigenvalues and speed: 1 pt.
\item Comparison with Solow: 0.5 pt.
}
\end{sol}

\textbf{c)} (4 points) Find the saving rates $(s_k,s_h)$ that maximize steady-state consumption per
effective worker $c^*=(1-s_k-s_h)y^*$, and prove that they are a global maximum. What happens to the
convergence speed and to $y^*$ as $\alpha+\beta\to1$, and what is the case $\beta\to0$?
\resolespaco
\begin{sol}
\passo{objective}{$\ln c^*=\ln(1-s_k-s_h)+\dfrac{\alpha\ln s_k+\beta\ln s_h}{1-\alpha-\beta}-\dfrac{\alpha+\beta}{1-\alpha-\beta}\ln N$.}
\passo{FOCs}{$\dfrac{-1}{1-s_k-s_h}+\dfrac{\alpha}{(1-\alpha-\beta)s_k}=0$ and
$\dfrac{-1}{1-s_k-s_h}+\dfrac{\beta}{(1-\alpha-\beta)s_h}=0$. Dividing, $s_k/s_h=\alpha/\beta$. Then
$(1-\alpha-\beta)s_k=\alpha(1-s_k-\tfrac\beta\alpha s_k)=\alpha-(\alpha+\beta)s_k$, so $s_k=\alpha$ and $s_h=\beta$.}
\passo{global maximum}{$\ln(1-s_k-s_h)$ is concave (log of an affine function) and
$\alpha\ln s_k+\beta\ln s_h$ is concave, so $\ln c^*$ is concave on the open simplex and the FOC point is
its global maximum. Equivalently, at $s_k=\alpha$ the net marginal product equals $N$:
$f_k=\alpha y^*/k^*=\alpha N/s_k=N$.}
\passo{limits}{As $\alpha+\beta\to1$: $\lambda=(1-\alpha-\beta)N\to0$ and the exponent $1/(1-\alpha-\beta)\to\infty$; the
model becomes AK, where saving rates change growth rates, not levels. As $\beta\to0$: the Solow
model, with $s_{gr}=\alpha$.}
\refbib{Romer (2012), problem 1.5; Kurlat (2020), \S4.3, eq.\ (4.3.4), pp.~61--62.}
\rubrica{
\item Objective in logs and FOCs: 1.5 pts.
\item Solution $s_k=\alpha$, $s_h=\beta$: 1 pt.
\item Global-maximum argument: 1 pt. Only FOCs, no concavity: 0 for this line.
\item Limits: 0.5 pt.
}
\end{sol}

\textbf{d)} (4 points) Use $\alpha=\beta=1/3$; Brazil: $s_k=0.18$, $s_h=0.10$, $N=0.065$; US: $s_k=0.22$,
$s_h=0.12$, $N=0.06$. Compute the predicted ratio $y^*_{BR}/y^*_{US}$. Output per worker is
$A\cdot y$; if the observed ratio of output per worker is $0.25$, compute the ratio of $A$ and of
Hicks-neutral TFP $B\equiv A^{1-\alpha-\beta}$. Split the log gap into accumulation and technology, and
say what the model predicts for the gap's dynamics.
\resolespaco
\begin{sol}
\passo{predicted ratio}{With $\alpha=\beta=1/3$ the coefficients in $\ln y^*$ are $1,1,-2$:
$\ln\frac{0.18}{0.22}+\ln\frac{0.10}{0.12}-2\ln\frac{0.065}{0.06}=-0.2007-0.1823-0.1601=-0.5431$, ratio $0.581$.}
\passo{technology}{$\dfrac{A_{BR}}{A_{US}}=\dfrac{0.25}{0.581}=0.430$; Hicks-neutral
$B_{BR}/B_{US}=0.430^{1/3}=0.755$. The two numbers describe the same gap under different
conventions.}
\passo{split}{$\ln0.25=-1.386$: accumulation $-0.543$ (39\%), technology $-0.843$ (61\%).}
\passo{dynamics}{Conditional convergence only: each country approaches its own $y^*$ at 2\% a year
(half-life 35 years), but the 61\% due to $A$ never closes inside the model. Both countries save
far below the Golden Rule ($s_k+s_h=0.28<2/3$), so raising saving costs consumption during the
transition.}
\intuicao{Human capital doubles the leverage of saving rates on income, yet even this generous
model leaves most of Brazil's gap to technology and efficiency.}
\refbib{Romer (2012), \S4.1--4.2; Kurlat (2020), \S5.3--5.5, pp.~80--95.}
\rubrica{
\item Predicted ratio: 1 pt.
\item $A$ and $B$ ratios with the convention explained: 1 pt.
\item Log split: 1 pt.
\item Dynamics and Golden Rule reading: 1 pt. Claiming absolute convergence: 0 for this line.
}
\end{sol}

\end{document}
